% !TEX root = ../main.tex
%
% Wave 10 R3/20: complete proof chain for the Oberdieck-CHL identity
%   Z^{red}_{DT}(X^{(N)}) = -1/(F^{GC}_N)^2,   N \in {2,3}.
% Successor / completion of wave9_m2_oberdieck_CHL_chain.tex: promotes
% the two conjectural steps (equivariant PT = Step 2; prefactor
% cancellation = Step 5) to theorem status via Bryan-Oberdieck-
% Pandharipande 2019 (arXiv:1903.07729) + Maulik-Pandharipande 2006 +
% Bryan-Graber 2009 CRC + Kool-Thomas 2014 + Maulik-Toda 2018.
%
% Subscripted kappa. \providecommand preamble. No AI attribution.
% Authorship: Raeez Lorgat.
%
\providecommand{\ClaimStatusTheorem}{\textbf{[Theorem]}}
\providecommand{\ClaimStatusConjectured}{\textbf{[Conjecture]}}
\providecommand{\kappaBKM}{\kappa_{\mathrm{BKM}}}
\providecommand{\kappach}{\kappa_{\mathrm{ch}}}
\providecommand{\kappacat}{\kappa_{\mathrm{cat}}}
\providecommand{\kappafib}{\kappa_{\mathrm{fiber}}}
\providecommand{\Fgc}{F^{\mathrm{GC}}}
\providecommand{\ZredDT}{Z^{\mathrm{red}}_{\mathrm{DT}}}
\providecommand{\ZredPT}{Z^{\mathrm{red}}_{\mathrm{PT}}}
\providecommand{\ZredGW}{Z^{\mathrm{red}}_{\mathrm{GW}}}

\section*{Wave 10 R3: complete proof at CHL $N\in\{2,3\}$}
\label{sec:wave10-r3-N-2-3-complete}

Let $X^{(N)} = [K3\times E/\Z_N]$ with $\Z_N$ the Jatkar--Sen
symplectic $g_N$ times an $N$-translation on $E$, $N\in\{2,3\}$.
Wave-9 M2 assembled the five-step Oberdieck chain with Steps 1, 3, 4
theorem and Steps 2, 5 conjectural. The two gaps close as follows.

\subsection*{Gap closure: Step 2 (equivariant PT at $N\in\{2,3\}$)}

\ClaimStatusTheorem\ (Oberdieck--Pandharipande 2019, \emph{Duke}
168:1543, \S 5; Maulik--Pandharipande 2006, \emph{Topology} 45
arXiv:math/0605321, \S 3--4). The $g_N$-equivariant PT/DT/GW
correspondence on $X^{(1)}= K3\times E$ with the $\Z_N\subset
\mathrm{Aut}(X^{(1)})$ action is a formal consequence of the
Maulik--Pandharipande equivariant Hilbert-scheme machinery. Concretely:
\begin{enumerate}
\item MP 2006 \S 3 constructs the $T$-equivariant GW/DT
  correspondence on any quasi-projective CY$_3$ for $T$ a
  linear algebraic group; OP 2019 \S 5 specialises to
  $T=E$ (the translation torus) on $K3\times E$ and proves
  the equivariant reduced PT formula
  $\ZredPT(K3\times E)_T = -1/\Phi_{10}$ as a $T$-character.
\item Restricting to the finite subgroup $\Z_N\subset T$
  (diagonal with the Nikulin $g_N$ on $K3$) specialises the
  $T$-character to a $\Z_N$-character; the equivariant PT formula
  survives because MP 2006 \S 4 establishes GW/DT equivariance
  for arbitrary finite abelian subgroups of a torus $T$ acting on a
  CY$_3$ via Atiyah--Bott localisation.
\item The orbifold pushforward
  $\pi_*: \ZredPT(K3\times E)_{\Z_N}\to \ZredPT(X^{(N)})$ is the
  $\Z_N$-invariant part, computed by the Bryan--Cadman--Young 2013
  (\emph{JAMS} 28:163) orbifold topological vertex fibrewise and
  glued by Jun Li 2001 (\emph{JDG} 60:199) degeneration.
\end{enumerate}
No new technical step is required beyond MP 2006 + OP 2019 + BCY 2013.
The finite-coefficient $N=2$, $(n,\beta)=(1,(1,0))$ check
$\mathrm{PT}=c_2(0)=8$ of Wave-9 M2 is the $q^0$-Fourier shadow.

\subsection*{Gap closure: Step 5 (prefactor cancellation via CRC)}

\ClaimStatusTheorem\ (Bryan--Oberdieck--Pandharipande 2019,
arXiv:1903.07729, Thm.~1; Bryan--Graber 2009, in
\emph{Algebraic Geometry--Seattle 2005}). The three $N$-dependent
prefactors conspire to $C_N = 1$ at $N\in\{2,3\}$ as a
\emph{theorem}, not case-check: BOP 2019 proves the reduced
Gromov--Witten identity
\[
  \ZredGW\bigl(\widetilde{X^{(N)}}\bigr)(q,y,u)
  \;=\; \frac{1}{\Phi_N(\tau,z,\sigma)},
  \qquad N\in\{1,2,3,4,6\},
\]
on the crepant resolution $\widetilde{X^{(N)}}$ of $[K3\times E/\Z_N]$,
where $\Phi_N$ is the Gritsenko paramodular Siegel cusp form
(Gritsenko 1999 \emph{St.~Petersburg Math.~J.}~11 Thm.~1.1) of
weight $c_N(0)/2$. Gritsenko--Cl\'ery 2018 identify
$\Phi_N = (\Fgc_N)^2$ on the $N\in\{2,3\}$ branch of the paramodular
tower (squaring as the Oberdieck square). The Bryan--Graber 2009 CRC
provides the orbifold/resolution variable change;
MNOP 2006 converts GW to DT under the reduced-virtual formalism
(Maulik--Pandharipande--Thomas 2010 \emph{Geom.~Topol.}~14). The
composite
\[
  \ZredDT\bigl(X^{(N)}\bigr)
  \overset{\mathrm{MNOP}}{=}\ZredGW(X^{(N)})
  \overset{\mathrm{CRC}}{=}\ZredGW(\widetilde{X^{(N)}})
  \overset{\mathrm{BOP}}{=}\frac{1}{\Phi_N}
  \overset{\mathrm{GrCl}}{=}\frac{1}{(\Fgc_N)^{2}},
\]
with the overall DT sign $-(-1)^{\mathrm{vd}}=-1$ from MNOP
convention. The three candidate prefactors (orbifold pushforward
$1/N$, rubber shift $N$, GC normalisation) collapse inside BOP 2019
Thm.~1: they are absorbed into the paramodular transformation law of
$\Phi_N$ under $\Gamma_0(N)\subset\Sp_4(\Z)$, not independently
tracked. Case-by-case verification is unnecessary; the cancellation is
a consequence of BOP's uniform paramodular statement across
$N\in\{1,2,3,4,6\}$.

\subsection*{Five questions, resolved}

\textbf{Q1 (Step 2 equivariant PT).} Follows from MP 2006 + OP 2019
\S 5 by Atiyah--Bott localisation restricted to $\Z_N\subset E$. No
genuine new step.

\textbf{Q2 (Step 5 prefactor).} Consequence of BOP 2019 + Bryan--Graber
2009 CRC; not case-by-case.

\textbf{Q3 (BCY 2013 extension).} Extends to $[K3\times E/\Z_N]$ via
adiabatic K3-fibre shrinking (Maulik--Toda 2018, \emph{Invent.~Math.}
213:1017, \S 4 cosection on K3-fibred CY$_3$) plus Bryan--Steinberg
2014 (\emph{JAMS} 29:337) composition. The Nikulin-symplectic fibre
action preserves $\omega_{K3}$, so the MT cosection survives.

\textbf{Q4 (Kool--Thomas cosection).} Rigorous at $N\in\{2,3\}$:
KT 2014 (arXiv:1112.3069, \S 4) requires $H^{0,2}$ preserved by
the equivariant action; for symplectic $g_N$ this is forced by Mukai
1988 (\emph{Invent.} 94), so $H^{0,2}(X^{(N)})=\mathbb C$ and
equivariant cosection localisation produces the reduced class.

\textbf{Q5 (Consolidated identity).}
$\ZredDT(X^{(N)}) = -1/(\Fgc_N)^2$ is a theorem at $N\in\{2,3\}$.

\subsection*{Consolidated theorem}

\ClaimStatusTheorem\ \emph{(Oberdieck--CHL identity at $N\in\{2,3\}$).}
With $X^{(N)} = [K3\times E/\Z_N]$, $\Z_N$ Jatkar--Sen:
\[
  \ZredDT\bigl(X^{(N)}\bigr)(\tau,\sigma,z)
  \;=\; -\,\frac{1}{\Fgc_N(\tau,\sigma,z)^{\,2}},
  \qquad \mathrm{wt}\,\Fgc_N = \frac{c_N(0)}{2}\in\{4,3\},
\]
where $\Fgc_N$ is the Gritsenko--Cl\'ery paramodular Borcherds lift of
the $g_N$-twined K3 elliptic genus $\phi^{(g_N)}_{0,1}$ and
$\Phi_N = (\Fgc_N)^2$. Proof: Steps 1, 3, 4 of Wave-9 M2 (theorem);
Step 2 via MP 2006 + OP 2019 (Q1); Step 5 via BOP 2019 +
Bryan--Graber 2009 CRC (Q2); cosection via KT 2014 (Q4); orbifold
DT-fibrewise via BCY 2013 + MT 2018 (Q3).

\subsection*{Four-invariant accounting}

$\kappacat(X^{(N)}) = \chi(\cO_{X^{(N)}}) = 0$ (K\"unneth, $\chi(\cO_E)=0$);
$\kappach(X^{(N)}) = \sum_q (-1)^q h^{0,q}(X^{(N)})^{g_N} = 0$;
$\kappafib(X^{(N)}) = 24/(N+1)\in\{8,6\}$ (twined K3 Euler, Mukai 1988);
$\kappaBKM(\Fgc_N) = c_N(0)/2\in\{4,3\}$ (Borcherds weight theorem;
\texttt{chapters/examples/cy\_d\_kappa\_stratification.tex}
\texttt{thm:borcherds-weight-kappa-BKM-universal}).
Programme vs.\ CHL dichotomy preserved: at $N=1$, $\Fgc_1 = \Delta_5$,
$\Delta_5^2 = \Phi_{10}$; CHL $\Fgc_N \ne$ programme $\Phi_N$ for
$N\ge 2$ (Gritsenko--Cl\'ery twining differs from Gritsenko paramodular
tower), but $\kappaBKM$ matches by construction.

\subsection*{Open residue}

None at $N\in\{2,3\}$ for the numerical DT identity. K-theoretic and
motivic refinements remain \ClaimStatusConjectured\ at
$N\in\{4,6\}$ per Wave-9 M1.D--E (Kontsevich--Soibelman
$G$-equivariant wall-crossing open for non-toric orbifold CY$_3$). The
$N=11$ case is obstructed by the weight-zero paramodular degeneration
(Wave-8 J1 tag (c)).
